%% ma: L12-B02-0001
%% lop: 12
%% chuong: 1
%% bai: 2
%% dang: 12.2.1
%% loai: DS
%% muc: [NB, TH, TH, VD]
%% dap_an: "ĐSĐĐ"
%% nguon: "0. ĐỀ THAM KHẢO TN THPT 2025.pdf (D00), Phần II Câu 1"
%% kien_thuc: "Bài 2 (giá trị lớn nhất trên đoạn), Bài 1 (cực trị); đạo hàm hàm lượng giác (Lớp 11 Bài 32), phương trình lượng giác cơ bản (Lớp 11 Bài 4) — các bài trước"
%% trang_thai: cho-duyet
%% ly_do: "Dạng toán mới đề xuất 12.2.1, chờ giáo viên duyệt"
\begin{ex}
Cho hàm số $f(x)=2\cos x+x$.
\choiceTF
{\True $f(0)=2$; $f\!\left(\dfrac{\pi}{2}\right)=\dfrac{\pi}{2}$.}
{Đạo hàm của hàm số đã cho là $f'(x)=2\sin x+1$.}
{\True Nghiệm của phương trình $f'(x)=0$ trên đoạn $\left[0;\dfrac{\pi}{2}\right]$ là $\dfrac{\pi}{6}$.}
{\True Giá trị lớn nhất của $f(x)$ trên đoạn $\left[0;\dfrac{\pi}{2}\right]$ là $\dfrac{\pi}{6}+\sqrt3$.}
\loigiai{a) $f(0)=2\cos0+0=2$; $f\!\left(\frac{\pi}{2}\right)=2\cos\frac{\pi}{2}+\frac{\pi}{2}=\frac{\pi}{2}$. Đúng.
b) $f'(x)=-2\sin x+1$. Sai.
c) $f'(x)=0\Leftrightarrow\sin x=\frac12$; trên $\left[0;\frac{\pi}{2}\right]$ chỉ có $x=\frac{\pi}{6}$. Đúng.
d) $f(0)=2$, $f\!\left(\frac{\pi}{6}\right)=\sqrt3+\frac{\pi}{6}\approx2{,}256$, $f\!\left(\frac{\pi}{2}\right)=\frac{\pi}{2}\approx1{,}571$. Giá trị lớn nhất là $\frac{\pi}{6}+\sqrt3$. Đúng.}
\end{ex}
